\scanpage{64}
\mspage{20}
Mais ces opérations de bouchage sont relativement superficielles, en cela qu'elles ne permettent pas de relier les $\mathcal T_{g,\nu}$, pour $g$ fixé, à des $\mathcal T_{g',\nu'}$ avec $g'<g$. Pour ceci, il faut l'opération \textcircled{$5^\circ$} de gommage des cercles de découpage (= ex-cercles de recollement !).

Néanmoins, il y a, \deleted{une} \inserted{en} liaison avec ça, une autre structure importante à expliciter, c'est celle \deleted{\unsure{des}} des groupoïdes fondamentaux \inserted{$\Pi_1(X;R)$} des $(X,[\unsure{\Sigma}],R)$, \inserted{pour $R\ne\varnothing$,} à base des $R$, (en tant que groupoïdes dépendant fonctoriellement bien sûr de $R$, et \deleted{la donnée} \unsure{comme} \unsure{nouvelle opération}) pour \unsure{tte} \deleted{\unsure{reconstitution}} \inserted{\unsure{représentation}} \unsure{du groupoïde par ens finis}, \unsure{d'un} \inserted{[\unsure{définissant un}]} revêtement \inserted{$X'$} de $X$), d'un nouvel objet du \deleted{$\unsure{\mathcal T_{\mathrm{top}}}$} groupoïde $S\mathcal T_{\mathrm{top}}^\pm$. Déjà dans le cas de $S_1\mathcal T_{0,3}$ la question se pose et est non triviale --- en fait, c'est sans doute là une exigence trop ambitieuse, car

\margin{$\Sigma$ ici est formellement \deleted{\unsure{vide}} vide, \unsure{autant donc} \unsure{pour les mettre dans les} \unsure{cas}, i.e. travailler dans $S\mathcal T_{\mathrm{top}}^\pm$ au lieu de $S\mathcal TD_{\mathrm{top}}^\pm$.}
\mspage{63}
